% ===== tr-05 块1c：第 4 章 pp.36--38（原文页 36--38）=====
% [原文 p.36]
% [Tutor Note: 原文把对数空间态射的底空间映射记作 \underline{F}（带下划线），把幺半群同态记作 F^\#（带上标 \#），二者不可混用；α_X∘F^\# = \underline{F}^*∘α_Y 中的星号属于底空间映射 \underline{F}。]
对数空间之间的态射 $F:X^\dagger\to Y^\dagger$ 由底空间之间的态射 $\underline{F}:X\to Y$ 以及一个同态 $F^\#:\underline{F}^{-1}(\mathcal{M}_Y)\to\mathcal{M}_X$ 组成，后者与结构同态交换：
\[
\alpha_X\circ F^\#=\underline{F}^*\circ\alpha_Y.
\]

关键的例子如下。

\begin{example}{4.8.2}
(1) 设 $X$ 是概形，$Y\subseteq X$ 是余维一的闭子集。记 $j:X\setminus Y\to X$ 为包含映射。则「在 $Y$ 之外可逆的正则函数层」的包含
\[
\alpha_X:\mathcal{M}_X=j_*(\mathcal{O}_{X\setminus Y}^\times)\cap\mathcal{O}_X\to\mathcal{O}_X
\]
是 $X$ 上的一个对数结构。它称为 $X$ 上的\emph{除子型对数结构}。

(2) 一个\emph{预对数结构}，即幺半群层的一个任意同态 $\varphi:\mathcal{P}\to\mathcal{O}_X$，通过
\[
\mathcal{M}_X=(\mathcal{P}\oplus\mathcal{O}_X^\times)/\{(p,\varphi(p)^{-1})\mid p\in\varphi^{-1}(\mathcal{O}_X^\times)\}
\]
定义一个相伴的对数结构 $\mathcal{M}_X$，且 $\alpha_X(p,h)=h\cdot\varphi(p)$。

(3) 若 $f:X\to Y$ 是概形的态射，而 $\alpha_Y:\mathcal{M}_Y\to\mathcal{O}_Y$ 是 $Y$ 上的对数结构，则由 $\alpha_Y:f^{-1}(\mathcal{M}_Y)\to f^{-1}\mathcal{O}_Y$ 与 $f^*:f^{-1}\mathcal{O}_Y\to\mathcal{O}_X$ 的复合给出的预对数结构 $f^{-1}(\mathcal{M}_Y)\to\mathcal{O}_X$ 在 $X$ 上定义一个相伴的对数结构，即\emph{拉回对数结构}。

(4) 在 (1) 中，我们可以用 (3) 把 $X$ 上的对数结构拉回到 $Y$。因此特别地，若 $\mathcal{X}\to D$ 是环面退化，则包含 $\mathcal{X}_0\subseteq\mathcal{X}$ 给出 $\mathcal{X}$ 上的一个对数结构，以及 $\mathcal{X}_0$ 上一个诱导的对数结构。类似地，包含 $0\in D$ 给出 $D$ 上的一个对数结构，以及 $0$ 上一个诱导的对数结构。这里 $\mathcal{M}_0=\mathbb{C}^\times\oplus\mathbb{N}$，其中 $\mathbb{N}$ 是自然数（非负整数）的（加法）幺半群，且
\[
\alpha_0(h,n)=\begin{cases} h & n=0,\\ 0 & n\neq 0.\end{cases}
\]
$0^\dagger$ 通常称为\emph{标准对数点}。于是我们有对数态射 $\mathcal{X}^\dagger\to D^\dagger$ 与 $\mathcal{X}_0^\dagger\to 0^\dagger$。
% [Tutor Note: $\mathcal{M}_0=\mathbb{C}^\times\oplus\mathbb{N}$ 中的 $\mathbb{N}$ 是\emph{加法}幺半群（非负整数），与 $\mathbb{C}^\times$ 的乘法结构作直积。$\alpha_0$ 是到乘法幺半群 $(\mathbb{C},\cdot)$ 的幺半群同态：$n=0$ 时 $\alpha_0(h,0)=h$，$n\neq 0$ 时取 $0$（$0$ 是乘法下的吸收元），故 $\mathbb{N}$ 的非零部分被整体压到 $0$。]

(5) 若 $\sigma\subseteq M_{\mathbb{R}}=\mathbb{R}^n$ 是严格凸有理多面体锥，$\sigma^\vee\subseteq N_{\mathbb{R}}$ 是对偶锥，令 $P=\sigma^\vee\cap N$：这是加法下的一个幺半群。由 $\sigma$ 定义的仿射环面簇可写成 $X=\operatorname{Spec}\mathbb{C}[P]$。于是我们有一个由幺半群同态
\[
P\to\mathbb{C}[P]
\]
（由 $p\mapsto z^p$ 给出）诱导的预对数结构。因此 $X$ 上有一个相伴的对数结构。它事实上与由 $\partial X\subseteq X$ 诱导的对数结构相同，其中 $\partial X$ 是 $X$ 的环面边界，即 $X$ 的所有环面除子之并。
% [原文 p.37]
% [Tutor Note: 「精细」（fine）的定义要求 $P\to P^{\mathrm{gp}}$ 是\emph{单射}，「饱和」（saturated）要求 $P$ 是整的且对 $m>0$、$mp\in P$ 有 $p\in P$；二者是不同条件，原文特意分开说明，翻译时不可合并成一句。]
若 $p\in P$，则单项式 $z^p$ 定义一个映射 $f:X\to\operatorname{Spec}\mathbb{C}[\mathbb{N}]=\mathbb{A}^1$，它关于 $\operatorname{Spec}\mathbb{C}[\mathbb{N}]$ 上由 $\mathbb{N}\to\mathbb{C}[\mathbb{N}]$ 类似地诱导的对数结构是一个对数态射。纤维 $X_0=\operatorname{Spec}\mathbb{C}[P]/(z^p)$ 是 $X$ 的子概形，$X_0$ 上有诱导的对数结构，并有如 (4) 中的映射 $X_0^\dagger\to 0^\dagger$。对数态射 $f$ 是\emph{对数光滑}态射的一个例子，见定义 4.8.3。

定义 4.7.1 的条件 (4) 事实上蕴含：在 $Z$ 之外局部地，$\mathcal{X}^\dagger$ 与 $\mathcal{X}_0^\dagger$ 都具有上述形式。因此我们应把 $\mathcal{X}^\dagger\to D^\dagger$ 视为在 $Z$ 之外对数光滑的；从对数的观点看，$\mathcal{X}_0^\dagger$ 在 $Z$ 之外可以当作非奇异概形一样处理。

(6) 给定如 (5) 中的幺半群 $P$ 以及态射 $X\to\operatorname{Spec}\mathbb{C}[P]$，我们可以把上面在 $\operatorname{Spec}\mathbb{C}[P]$ 上定义的对数结构拉回到 $X$。若一个对数概形 $X^\dagger$ 在 étale 局部上可以这样描述，我们就说 $X^\dagger$ 是\emph{精细饱和}（fine saturated）对数概形。形容词「精细」告诉我们，它在局部上通过映入形如 $\operatorname{Spec}\mathbb{C}[P]$ 的概形来描述，其中 $P$ 是有限生成\emph{整}幺半群，即典范同态 $P\to P^{\mathrm{gp}}$ 是单射。形容词「饱和」告诉我们幺半群 $P$ 是饱和的；这指的是 $P$ 是整的，且只要 $p\in P^{\mathrm{gp}}$ 对某个 $m>0$ 满足 $mp\in P$，就有 $p\in P$。这样的幺半群例如作为有理多面体锥与格的交而出现。

关于对数几何的文献大多只适用于精细对数结构。在关键例子 $\mathcal{X}_0^\dagger\to 0^\dagger$ 中，对数结构在集合 $Z$ 之外是精细饱和的。然而一般地它在 $Z$ 上并不精细，这往往造成许多技术困难，因为必须发展新的技术才能正确处理 $Z$ 上的对数结构。\proofend
\end{example}

对数光滑的概念推广了例 4.8.2 的 (5) 中的那些态射：

\begin{definition}{4.8.3}
精细对数概形的态射 $f:X^\dagger\to Y^\dagger$ 称为\emph{对数光滑}的，如果它在 $X$ 与 $Y$ 上 étale 局部地可放进一个交换图表
\begin{center}
\begin{tikzcd}
X \arrow[r] \arrow[d] & \operatorname{Spec}\mathbb{Z}[P] \arrow[d] \\
Y \arrow[r] & \operatorname{Spec}\mathbb{Z}[Q]
\end{tikzcd}
\end{center}
并具有以下性质：
\begin{enumerate}
\item 例 4.8.2 的 (5) 中 $\operatorname{Spec}\mathbb{Z}[P]$ 与 $\operatorname{Spec}\mathbb{Z}[Q]$ 上的典范对数结构，分别拉回为 $X$ 与 $Y$ 上的对数结构。
\item 诱导出的态射
\[
X\to Y\times_{\operatorname{Spec}\mathbb{Z}[Q]}\operatorname{Spec}\mathbb{Z}[P]
\]
是概形的光滑态射。
\item 右方的竖箭头由幺半群同态 $Q\to P$ 诱导，且 $\ker(Q^{\mathrm{gp}}\to P^{\mathrm{gp}})$ 与 $\operatorname{coker}(Q^{\mathrm{gp}}\to P^{\mathrm{gp}})$ 的挠部分都是阶在 $X$ 上可逆的有限群。这里 $P^{\mathrm{gp}}$ 表示 $P$ 的 Grothendieck 群。
\end{enumerate}
\end{definition}
% [原文 p.38]
对数光滑态射在最简单的情形下包含正规交叉态射。

对数概形 $X^\dagger$ 上总有一个正合列
\[
1\longrightarrow\mathcal{O}_X^\times\xrightarrow{\alpha^{-1}}\mathcal{M}_X\longrightarrow\overline{\mathcal{M}}_X\longrightarrow 0,
\]
其中我们把幺半群的商层 $\overline{\mathcal{M}}_X$ 按加法书写。我们称 $\overline{\mathcal{M}}_X$ 为该对数结构的\emph{幽灵层}（ghost sheaf）。
% [Tutor Note: 正合列 1→O_X^× →(α^{-1}) M_X → \overline{M}_X → 0 中，第二个箭头是 α_X 在 O_X^× 上的限制之\emph{逆}（α_X 在单位群上同构），不是 α_X 本身；抄写时容易漏掉指数 −1。]
我喜欢把 $\overline{\mathcal{M}}_X$ 看作指定了与该对数结构相伴的组合信息。例如，若 $X^\dagger$ 由 Cartier 除子 $Y\subseteq X$ 诱导（$X$ 正规），则 $x\in X$ 处的茎 $\overline{\mathcal{M}}_{X,x}$ 就是 $x$ 的某个邻域上支集含于 $Y$ 的有效 Cartier 除子构成的幺半群。

为理解对数结构的拉回，有用的一点是：若 $f:Y\to X$ 是态射，$X$ 带有对数结构，且 $Y$ 被赋予拉回对数结构，则 $\overline{\mathcal{M}}_Y=f^{-1}\overline{\mathcal{M}}_X$。在 $\mathcal{M}_X$ 由包含 $Y\subseteq X$ 诱导的情形，$\overline{\mathcal{M}}_X$ 的支集在 $Y$ 上，于是我们可以把 $\overline{\mathcal{M}}_X$ 与 $\overline{\mathcal{M}}_Y$ 等同起来，前者是 $X$ 上除子型对数结构的幽灵层，后者是它限制到 $Y$ 上的幽灵层。

\begin{exercise}{4.8.4}
证明：在例 4.8.2 的 (5) 中，若 $\dim\sigma=\dim M_{\mathbb{R}}$ 且 $x$ 是 $X$ 的唯一零维环面轨道，则 $\overline{\mathcal{M}}_{X,x}=P$。更一般地，当 $x\in X$ 位于对应于 $\sigma$ 的一个面 $\tau$ 的环面轨道中时，
\[
\overline{\mathcal{M}}_{X,x}=\frac{\tau^\vee\cap N}{\tau^\perp\cap N}=\operatorname{Hom}_{\mathrm{monoid}}(\tau\cap M,N).
\]
特别地，$\tau$ 可以由 $\operatorname{Hom}_{\mathrm{monoid}}(\overline{\mathcal{M}}_{X,x},\mathbb{R}_{\ge 0})$ 还原出来，其中 $\mathbb{R}_{\ge 0}$ 是非负实数加法幺半群。\proofend
\end{exercise}

在接下来的几节中，我们所考虑的关键对数空间将是由环面退化 $\mathcal{X}\to D$ 产生的那些。如上所述，中心纤维 $\mathcal{X}_0\subseteq\mathcal{X}$ 在 $\mathcal{X}$ 上诱导一个除子型对数结构，限制之后给出一个对数概形 $\mathcal{X}_0^\dagger$，连同态射 $\mathcal{X}_0^\dagger\to 0^\dagger$，后者在坏集 $Z\subseteq\mathcal{X}_0$ 之外是对数光滑的。

现在我们可以用下面的图表来进一步阐述我们希望采取的观点：
\begin{center}
\begin{tikzpicture}[node distance=1.0cm]
\node[draw,ellipse,minimum width=2.6cm,minimum height=1.05cm] (A) {A-模型};
\node[draw,ellipse,right=of A,minimum width=3.3cm,minimum height=1.05cm] (L) {对数几何};
\node[draw,ellipse,right=of L,minimum width=2.6cm,minimum height=1.05cm] (B) {B-模型};
\draw[<->] (A) -- (L);
\draw[<->] (L) -- (B);
\end{tikzpicture}
\end{center}

镜像对称有两面。A-模型这一面涉及计数曲线：我们希望数一个环面退化 $\mathcal{X}\to D$ 的一般纤维中的曲线。有充分理由相信，这一计数实际上可以在 $\mathcal{X}_0^\dagger$ 上完成，用的是对数 Gromov--Witten 不变量的理论：见第 9 节。希望在于 $\mathcal{X}_0$ 是一个足够组合化的对象，以致这样的计数可以以组合方式完成。

$B$-面涉及复结构的形变。其想法是：要理解复结构的形变，我们应当从中心纤维 $\mathcal{X}_0^\dagger$ 出发，尝试构造光滑化，即构造以此为中心纤维的环面退化。对数结构对于找到唯一的光滑化是必要的。
